% !TeX root = ../main.tex
\section{Introduction}
\label{sec:introduction}

We consider smooth nonconvex--strongly-concave (NC--SC) minimax
optimization
\begin{equation}
    \min_{x\in\mathbb{R}^{d_x}}
    \max_{y\in\mathbb{R}^{d_y}} f(x,y),
    \label{eq:minimax-intro}
\end{equation}
where $f(x,\cdot)$ is strongly concave for every $x$, while
$f(\cdot,y)$ may be nonconvex. Strong concavity makes the inner
maximizer
\[
    y^*(x):=\arg\max_y f(x,y)
\]
well defined and gives rise to the primal function
\[
    P(x):=\max_y f(x,y).
\]
Under standard smoothness assumptions, $P$ is twice differentiable
but generally nonconvex. An approximate first-order stationary point
of $P$ therefore need not capture local optimality, which motivates
the search for approximate second-order stationary points.

Second-order algorithms for NC--SC minimax optimization have been
developed using cubic regularization, trust-region methods, and other
Newton-type updates
\citep{luo2022finding,chen2023cubic,yao2025trust,
wang2025gradient,chen2026hsda}.
A representative method is Minimax Cubic Newton (MCN)
\citep{luo2022finding}, which approximates the gradient and Hessian of
$P$ through inner maximization and achieves the characteristic
$O(\epsilon^{-3/2})$ second-order oracle complexity. Several subsequent
methods attain the same $\epsilon^{-3/2}$-type outer dependence under
related second-order stationarity criteria. More recently,
joint-minimization reformulations have led to methods that avoid
repeatedly solving the inner maximization
\citep{ma2026linesearch,ma2026singleloop}.

For methods that construct approximate derivatives of $P$ through
inner maximization, however, the outer iteration complexity alone does
not describe the total computational work. At each outer iterate
$x_k$, one must approximate $y^*(x_k)$ accurately enough that the
resulting first- and second-order information of $P$ can be used
reliably. In MCN, for example, accelerated gradient ascent is used for
this purpose. Its first-order oracle complexity is stated in
$\widetilde O(\epsilon^{-3/2})$ form, and the explicit bound contains a
logarithmic dependence on the target accuracy multiplying the leading
outer-order term. This raises a natural question: can one retain the
$O(\epsilon^{-3/2})$ second-order iteration dependence while removing
the multiplicative logarithmic factor from the total inner-gradient
complexity?

We answer this question affirmatively by combining existing
primal-function estimators with an analysis of the accumulated inner
work. The implicit gradient estimator studied by
\citet{ablin2020superefficiency} has known quadratic accuracy with
respect to the inner solution error; see their Proposition~2.3.
The same total-gradient and Schur-complement expressions also appear
in Newton-type minimax methods \citep{zhang2020newton}.
We use these quantities, together with a Newton correction of the
value, as oracle ingredients and derive the explicit error bounds
needed under our NC--SC assumptions.
The complexity improvement comes from combining these bounds with
warm starts and a cumulative cubic bound on the outer steps.
This path estimate controls the total cost of restoring inner
accuracy across successive maximization problems.

Our outer constructions build on the universal trust-region (UTR)
framework of \citet{jiang2026utr}. For the approximate
primal-function oracles required here, we also establish a
certified-inexact UTR complexity result with additive search overhead;
the generic result is deferred to the appendix so that the main text
can focus on the minimax problem. In the known-parameter setting, we
combine a fixed-scale UTR update with an abstract residual-based inner
solver. The inner-work analysis applies to any solver satisfying an
accelerated residual-reduction bound, with accelerated gradient descent
providing a concrete realization. The resulting method has
$O(\epsilon^{-3/2})$ leading total inner-gradient dependence, without
a multiplicative $\log(1/\epsilon)$ factor.

We then remove prior knowledge of the problem regularity constants
from both the outer and inner components. An adaptive UTR construction
combined with the parameter-free SCAR solver yields a fully
parameter-free method, but requiring uniformly accurate inner solves
reintroduces a multiplicative $\log(1/\epsilon)$ factor in the total
inner-gradient complexity. Our final method instead uses a
gradient-scaled adaptive trust-region scheme together with persistent
SCAR calls across successive inner maximization problems. By retaining
the curvature information learned by the inner solver and adapting the
inner accuracy to the current outer scale, it recovers the
$O(\epsilon^{-3/2})$ leading dependence in both outer and total
inner-gradient complexity.


\subsection{Main Contributions}
\label{sec:contributions}

Our main contributions are summarized as follows.

\paragraph{Sharper total inner-gradient complexity.}
Let $\ell$ and $\rho$ denote the Lipschitz constants of the gradient
and Hessian of $f$, respectively, let $\mu$ be the strong concavity
constant, and set
\[
    \kappa=\ell/\mu,
    \qquad
    \Delta_P=P(x_0)-\inf_xP(x).
\]
Algorithm~1 combines the existing implicit gradient and
Schur-complement Hessian with a fixed-scale inexact UTR scheme and
accesses the inner maximization through an observable
gradient-residual condition. The key estimate is a cumulative cubic
bound on the outer steps, which amortizes the logarithmic costs of
successive warm-started inner calls. The analysis is stated for a
general accelerated residual solver; accelerated gradient descent is
one concrete realization. The number of outer trust-region subproblems has
leading complexity
\[
    O\!\left(
        \Delta_P\kappa^{3/2}\sqrt{\rho}\,
        \epsilon^{-3/2}
    \right),
\]
while the accelerated-gradient realization has leading total
inner-gradient complexity
\[
    O\!\left(
        (1+\log\kappa)\Delta_P\kappa^2\sqrt{\rho}\,
        \epsilon^{-3/2}
    \right).
\]
All dependence on $\log(1/\epsilon)$ in the full bound is additive
rather than multiplicative in the leading $\epsilon^{-3/2}$ term.
In particular, this sharpens the explicit first-order oracle bound of
MCN with respect to the target accuracy. To our knowledge, this is the
first explicit total inner-gradient bound for a second-order NC--SC
method that approximates derivatives of the primal function through
inner maximization whose leading $O(\epsilon^{-3/2})$ term is not
multiplied by a logarithmic factor in $1/\epsilon$.
As a technical ingredient, we also establish a certified-inexact UTR
complexity result with additive search overhead; its generic statement
and proof are deferred to the appendix.

\paragraph{Fully parameter-free second-order methods.}
We further remove the regularity constants from the executable
algorithms. Algorithm~2 combines an adaptive UTR outer scheme with
SCAR and requires no knowledge of $\ell$, $\mu$, or $\rho$. It retains
the leading $O(\epsilon^{-3/2})$ outer complexity, but uniformly
accurate SCAR solves introduce a multiplicative $\log(1/\epsilon)$
factor in the total inner-gradient complexity. Algorithm~3 replaces
these independent high-accuracy solves by persistent SCAR tracking
within a gradient-scaled adaptive trust-region scheme. For every fixed
set of regularity constants and all sufficiently small $\epsilon$, it
reaches an
$O(\epsilon,\sqrt{\kappa^3\rho\epsilon})$ second-order stationary
point using
\[
    O\!\left(
        \Delta_P\kappa^{3/2}\sqrt{\rho}\,
        \epsilon^{-3/2}
    \right)
\]
trust-region subproblems and
\[
    O\!\left(
        \Delta_P\kappa^2\sqrt{\rho}\,
        \epsilon^{-3/2}
    \right)
\]
inner-gradient evaluations in their leading terms. The remaining
dependence on $1/\epsilon$ is additive and polylogarithmic. The
executable algorithm uses only the target tolerance and universal
numerical constants. The guarantee is class-uniform: for fixed
regularity constants, one sufficiently small accuracy threshold
applies to every admissible objective and initialization, although
this threshold is not known to the algorithm.


% !TeX root = ../main.tex
\begin{table}[t]
\centering
\caption{
Complexity comparison of global deterministic second-order methods
for smooth NC--SC minimax optimization. Only complexity bounds
explicitly stated in the cited papers are reported. Dependence on
problem regularity constants is displayed when stated in a directly
comparable form; initial optimality gaps and numerical constants are
suppressed.
}
\label{tab:ncsc-complexity}

\small
\renewcommand{\arraystretch}{1.18}
\setlength{\tabcolsep}{3pt}

\begin{tabularx}{\textwidth}{@{}
    >{\raggedright\arraybackslash}X
    >{\centering\arraybackslash}p{0.24\textwidth}
    >{\centering\arraybackslash}p{0.31\textwidth}
    >{\centering\arraybackslash}p{0.11\textwidth}
@{}}
\toprule
Method
& Outer complexity
& \shortstack[c]{Total inner-gradient\\evaluations}
& \shortstack[c]{Parameter-\\free}
\\
\midrule

MCN \citep{luo2022finding}
&
$O\!\left(
\kappa^{3/2}\sqrt{\rho}\,
\epsilon^{-3/2}
\right)$
&
$\widetilde O\!\left(
\kappa^2\sqrt{\rho}\,
\epsilon^{-3/2}
\right)$
&
No
\\[1mm]

MINIMAX-TR / MINIMAX-TRACE
\newline \citep{yao2025trust}
&
$O(\epsilon^{-3/2})$
&
Not reported
&
No
\\[1mm]

GRTR / LMNegCur
\newline \citep{wang2025gradient}
&
$\widetilde O\!\left(
\kappa^{3/2}\sqrt{\rho}\,
\epsilon^{-3/2}
\right)$
&
$\widetilde O\!\left(
\kappa^2\sqrt{\rho}\,
\epsilon^{-3/2}
\right)$
&
No
\\[1mm]

HSDA / IHSDA
\newline \citep{chen2026hsda}
&
$\widetilde O(\epsilon^{-3/2})$
&
Not reported
&
No
\\

\midrule

\textbf{Algorithm~1 (ours)}
&
$O\!\left(
\kappa^{3/2}\sqrt{\rho}\,
\epsilon^{-3/2}
\right)$
&
$O\!\left(
(1+\log\kappa)\kappa^2\sqrt{\rho}\,
\epsilon^{-3/2}
\right)$
&
No
\\[1mm]

Algorithm~2 (ours)
&
$O\!\left(
\kappa^{3/2}\sqrt{\rho}\,
\epsilon^{-3/2}
\right)$
&
$O\!\left(
\kappa^2\sqrt{\rho}\,
\epsilon^{-3/2}\log(1/\epsilon)
\right)$
&
Yes
\\[1mm]

\textbf{Algorithm~3 (ours)}
&
$O\!\left(
\kappa^{3/2}\sqrt{\rho}\,
\epsilon^{-3/2}
\right)$
&
$O\!\left(
\kappa^2\sqrt{\rho}\,
\epsilon^{-3/2}
\right)$
&
\textbf{Yes}
\\

\bottomrule
\end{tabularx}

\vspace{1mm}
\begin{minipage}{0.97\textwidth}
\footnotesize
Here $\kappa=\ell/\mu$, and $\rho$ is the Hessian-Lipschitz constant
of $f$. The notation $\widetilde O(\cdot)$ is retained whenever it is
used in the corresponding reference; hidden logarithmic factors in a
cited result are not expanded or reconstructed here.
``Not reported'' means that the method uses inner maximization but the
cited work does not state a separate comparable bound on the total
number of inner-gradient evaluations.
For Algorithms~1--3, the table displays the leading terms of the
bounds proved in this paper. The complete bounds for Algorithms~1
and~3 also contain lower-order additive logarithmic or polylogarithmic
terms in $1/\epsilon$.
\end{minipage}
\end{table}


Table~\ref{tab:ncsc-complexity} is intentionally restricted to methods
whose stated second-order guarantees and oracle counts are directly
comparable to the quantities considered here.
Cubic-LocalMinimax \citep{chen2023cubic}, whose complexity is stated
using a different local-minimax stationarity measure, is discussed
below but is not reparameterized for the table.
Likewise, the recent single-loop regularized Newton method of
\citet{ma2026singleloop} does not use an inner maximization loop of the
form measured by the third column and is therefore not included in
that comparison.


\subsection{Related Work}

\paragraph{Implicit differentiation and primal-function estimators.}
\citet{ablin2020superefficiency} analyze the implicit gradient
estimator for functions defined through optimization and establish
its quadratic error in Proposition~2.3.
In minimax optimization, \citet{zhang2020newton} use the corresponding
total gradient and Schur-complement Hessian; see their
equations~(1.1)--(1.2).
These estimators are ingredients of the present work.
Our contributions concern their integration into inexact UTR methods,
the accumulated inner-gradient complexity, and parameter-free
tracking across changing inner objectives.

\paragraph{First-order methods for NC--SC minimax optimization.}
First-order methods for NC--SC minimax problems have been studied
extensively. Representative approaches include gradient
descent-ascent, multi-step ascent, proximal schemes, and accelerated
methods
\citep{nouiehed2019solving,lin2020gradient,lin2020near}.
Their complexity for finding first-order stationary points of the
primal function has also been studied from both upper- and lower-bound
perspectives; see, for example, \citet{zhang2021complexity} and the
references therein. These results form the first-order counterpart to
the second-order stationarity problem considered in this paper.

\paragraph{Second-order methods for NC--SC minimax optimization.}
\citet{luo2022finding} proposed MCN, which alternates accelerated
gradient ascent for approximating $y^*(x)$ with a cubic-regularized
update based on approximate first- and second-order information of the
primal function. They established
$O(\kappa^{3/2}\sqrt{\rho}\,\epsilon^{-3/2})$ second-order oracle
complexity and
$\widetilde O(\kappa^2\sqrt{\rho}\,\epsilon^{-3/2})$ first-order oracle
complexity.
\citet{chen2023cubic} proposed Cubic-LocalMinimax and developed
inexact and stochastic variants, using cubic regularization to obtain
global convergence toward local minimax points. Their complexity is
formulated using a different stationarity measure, so we do not
convert it to the accuracy parameterization used in
Table~\ref{tab:ncsc-complexity}.

Trust-region approaches were developed by \citet{yao2025trust}, who
proposed MINIMAX-TR and MINIMAX-TRACE.
\citet{wang2025gradient} subsequently proposed gradient-norm
regularized trust-region and Levenberg--Marquardt methods, together
with inexact variants, and established
$\widetilde O(\kappa^{3/2}\sqrt{\rho}\,\epsilon^{-3/2})$ outer
complexity and
$\widetilde O(\kappa^2\sqrt{\rho}\,\epsilon^{-3/2})$ total
gradient-ascent complexity.
More recently, \citet{chen2026hsda} proposed homogeneous second-order
descent-ascent methods and developed an inexact Lanczos implementation
for large-scale problems.

A recent alternative is to avoid repeated inner maximization
altogether. \citet{ma2026linesearch} introduced an equivalent
joint-minimization reformulation that preserves first- and
second-order stationarity and used it to develop a parameter-free
line-search framework. Building on this reformulation,
\citet{ma2026singleloop} proposed a single-loop regularized Newton
method with $O(\epsilon^{-3/2})$ global iteration complexity and local
superlinear convergence. This line of work addresses the cost of the
inner loop by changing the formulation, whereas the present paper
studies the total inner-gradient complexity when approximate
derivatives of the primal function are constructed through inner
maximization.

\paragraph{Adaptive and parameter-free optimization.}
Adaptive optimization methods aim to reduce or eliminate prior
knowledge of problem regularity constants. In smooth nonconvex
optimization, the universal trust-region (UTR) framework
\citep{jiang2026utr} uses gradient-scaled trust-region updates together
with adaptive scale adjustment, and provides the outer optimization
framework on which our methods are built. On the inner side,
\citet{lan2026optimal} developed SCAR, a parameter-free accelerated
method for smooth strongly convex optimization that adapts to unknown
smoothness and strong-convexity parameters while retaining accelerated
gradient complexity. These two ingredients play different roles in
our development: UTR provides the outer second-order mechanism, whereas
SCAR provides a parameter-free realization of the strongly concave
inner maximization.

Related adaptive and parameter-free methods have also been developed
for minimax optimization and monotone variational inequalities
\citep{jiang2024adaptive,jiang2025generalized}.
More recently, \citet{wang2026parameterfree} proposed a fully
parameter-free second-order method for convex--concave minimax
optimization, requiring neither the Hessian-Lipschitz constant nor an
a priori distance bound to a solution. For NC--SC minimax problems,
\citet{ma2026linesearch} developed a parameter-free line-search
framework based on a joint-minimization reformulation.
In contrast, our constructions retain the primal-function formulation:
Algorithm~2 combines adaptive UTR with SCAR, while Algorithm~3 couples
a gradient-scaled adaptive trust-region scheme with persistent SCAR
tracking across successive inner maximization problems.


\subsection{Organization}

The remainder of the paper is organized as follows.
Section~\ref{sec:preliminaries} presents the problem setting,
properties of the primal function, the inner residual interface, and
the Newton-corrected estimators.
Section~\ref{sec:known-parameter} develops Algorithm~1, a
known-parameter UTR method with a general accelerated inner residual
solver, and establishes its outer and total inner-gradient
complexity.
Section~\ref{sec:parameter-free} first combines adaptive UTR with SCAR
in Algorithm~2 and then develops the persistent-SCAR
gradient-scaled trust-region method, Algorithm~3.
Section~\ref{sec:experiments} presents the numerical experiments, and
Section~\ref{sec:conclusion} concludes the paper.
The generic certified-inexact UTR results and deferred technical
proofs are collected in the appendices.
